\documentclass[ekobook.tex]{subfiles}

\begin{document}\setWPP
\setlist[itemize]{
  itemsep=0pt,
  parsep=0pt,      % space between paragraph lines inside one item
  topsep=0.5ex,    % space before/after the whole list (usually keep a little)
  partopsep=0pt
}
\setlength{\columnsep}{1em}
\lsection{Podnikatel a banka}
\begin{multicols}{2}
\qaw{Kdo může podnikat}{%
Podnikat mohou jak \textbl{fyzické osoby}, \wemph{jednotlivě}, pod svým vlastním jménem, tak \textbl{právnické osoby}, coby \wemph{skupina osob} podnikajících pod společným jménem.
}
\qaw{Jakými zákony se řídí podnikání FO}{%
Podnikání fyzických osob se řídí \texthl{zákonem o živnostenském podnikání a živnostenských úřadech}.
}
\qaw{Kde musí být zapsány podnikající FO}{%
Všechny fyzické osoby podnikající na území České Republiky musí být povinně zapsány \texthl{v živnostenském rejstříku}.
}
\qaw{Definice živnosti}{%
Živnost je: \texthl{soustavná činnost}; provozovaná pod \texthl{vlastním jménem}; na \texthl{vlastní odpovědnost}; za \texthl{účelem dosažení zisků}; a za \texthl{podmínek stanovených zákonem}.
}
\qaw{Všeobecné podmínky pro získání ŽL}{%
\textbl{Dovršení plnoletosti} (dovršení 18 let)\\
\textbl{Svéprávnost} (na základě lékařské zprávy může odebrat soud)\\
\textbl{Bezúhonnost v oboru podnikání} (prokazuje se výpisem z~rejstříku trestů)\\
\textbl{Nesmí mít dluhy vůči státu}
}
\qaw{Zvláštní podmínky pro získání ŽL}{%
\textbl{Odborná způsobilost} (výuční list; maturita; vysoká škola)
}
\qaw{Uveď příklady ohlašovacích živností a kdy mohu začít podnikat}{%
Podnikatel může začít podnikat v \textbl{den nahlášení} na \textbl{Živnostenský úřad}.\\
\begin{itemize}[label={\color{WPP}\faSlackHash}]
\item \textbl{řemeslné} -- (truhlář, tesař, kadeřnice) další podmínkou je zde \texthl{vyučení v oboru}, či \texthl{6letá praxe}
\item \textbl{vázané} -- (provoz autoškoly, revize elektrospotřebičů) dokládá se \texthl{odborná způsob} (obvykle určitým \wemph{certifikátem}) a~praxe v~oboru
\item \textbl{volné} -- (velkoobchod, maloobchod) stačí pouze všeobecné podmínky
\end{itemize}
}
\qaw{Uveď příklad koncesované živnosti a kdy mohu začít podnikat}{%
Podnikatel může začít podnikat až v den, kdy mu přijde \textbl{koncesní listina}.
Takzvaná \wemph{koncese} je povolení od státu k výkonu živnosti po dokázání způsobilosti v oboru.
Patří sem například: \wemph{ostraha majetku osob}; \wemph{provádění pyrotechnického průzkumu}; \wemph{provozování střelnic}
}
\qaw{Zánik živnostenského podnikání}{%
Živnostenské podnikání může zaniknout z několika důvodů.
Nejčastějšími způsoby jsou \texthl{na žádost podnikatele o ukončení činnosti} a \texthl{uplynutí lhůty na kterou bylo zřízené}.
Dále končí \texthl{smrtí podnikatele} nebo na \texthl{rozhodnutí soudu}.
}
\end{multicols}
\wbreak
\qaw{Co musí předložit podnikatel při založení bankovního účtu}{%
Pro zřízení \texthl{podnikatelského účtu} musí podnikatel předložit svůj \wemph{občanský průkaz} a \wemph{živnostenský list}.
}
\qaw{Jaké druhy příkazů k úhradě používá podnikatel}{%
\begin{itemize}[label={\color{WPP}\faSlackHash}]
\item \textbl{jednorázový příkaz} -- využívaný k placení momántálně vzniklých faktur; vždy se liší doba a částka provedení
\item \textbl{hromadný příkaz} -- více plateb provedených v jeden moment
\item \textbl{trvalý příkaz k úhradě} -- využívaný k opakovaným platbám; zálohy na daně, nájem, spoření
\item \textbl{příkaz k inkasu} -- dám povolení poskytovateli služeb, aby si mohl strhávat platby dle míry, jak jsem jejich služby využíval daný měsíc (je nastaven limit pro maximální možnou platbu)
\end{itemize}
}
\qaw{Vysvětli rozdíl mezi kreditní a debetní kartou}{%
K vlastnímu \wemph{bankovnímu účtu} dostaneme \textbl{debetní kartu}, která nám usnadňuje manipulaci s našimi zdroji a dává nám možnost platit přímo skrz platební terminály.\\
K \wemph{úvěrovému účtu} dostaneme \textbl{kreditní kartu}, která splňuje podobné funkce, ale povoluje nám čerpat do záporných čísel, čímž nám vzniká úvěr vůči bance, který musíme splatit (bezúročně, jestli jej splatíme v rámci smluveného období).
}
\qaw{Vysvětli tyto pojmy:}{%
\textbl{bankovní tajemství}
-- banka nesmí sdělovat informace o účtu, včetně výše uložené částky či provedených transakcí\\
\textbl{dispoziční právo}
-- svěřené oprávnění nakládat s bankovním účtem (v manželství manželce, ve firmě účetní); oprávnění můžeme omezit jen do určitých částek a na určité operace\\
\textbl{pojištění vkladů ve výši}
-- naše vklady jsou pojištěné do výše 100 000 €\\
\textbl{cash back{\color{antiDP}$_{[1]}$}}
-- obchodníci s logem Cashback nabízejí výběr peněz z bankovního účtu skrz pokladnu; obchodník vám odečte peníze z karty a dá vám hotovost na ruku\\
\textbl{cash back{\color{antiDP}$_{[2]}$ (více známá služba se stejným názvem, zmiňuji jenom jako bonus)}}
-- vrácení peněz při nákupu z vybraných e--shopů, (například služba Tipli), výše vrácených peněz se pohybuje mezi 2 \% až 10 \%, e-shopům se vyplatí do tohoto systému vstoupit, protože to pro ně funguje jako reklama
}
\qaw{Jak se dělí termínované vklady podle času}{%
Finanční vklady do banky, které \texthl{nemůžeme vybrat před dobou splatnosti} (případně musíme \wemph{zaplatit sankce} za dřívější vybrání); banka k vkladu přičíta \texthl{pevně stanovený úrok}.\\
\textbl{krátkodobé} -- do jednoho roku\\
\textbl{střednědobé} -- mezi rokem až pěti\\
\textbl{dlouhodobé} -- nad pět let
}
\wpage
\end{document}